\documentclass[11pt]{article}
\usepackage[a4paper,margin=1.9cm]{geometry}
\usepackage[T1]{fontenc}
\usepackage{lmodern}
\usepackage{amsmath,amssymb,amsthm,mathtools}
\usepackage{graphicx}
\usepackage{xcolor}
\usepackage{microtype}
\usepackage{enumitem}
\usepackage[hidelinks]{hyperref}

\definecolor{ink}{HTML}{172033}
\definecolor{accent}{HTML}{5C3566}
\definecolor{soft}{HTML}{F6F0F8}
\color{ink}
\newtheorem{theorem}{Theorem}
\newtheorem{lemma}{Lemma}
\newcommand{\FBL}{\operatorname{FBL}_{\mathbb C}}
\newcommand{\R}{\mathbb R}
\newcommand{\C}{\mathbb C}
\newcommand{\epsf}{\varepsilon}

\title{Fourier evaluation proves spectral inclusion for free complex Banach lattices}
\author{Solution packet for arXiv:2207.08090}
\date{11 August 2026}

\begin{document}
\maketitle

\begin{center}
\fcolorbox{accent}{soft}{\parbox{0.91\textwidth}{
\textbf{Outcome.} For every complex Banach space $E$ and every bounded operator $T:E\to E$, the induced lattice homomorphism $\overline T:\FBL[E]\to\FBL[E]$ satisfies
\[
 \sigma(T)\subseteq\sigma(\overline T).
\]
The missing complex phase is recovered by the first Fourier coefficient of the circle of evaluation functionals associated with an eigenfunctional of $T^*$.}}
\end{center}

\section*{Source question}
The source constructs $\FBL[E]$ as a complexification of a real free Banach lattice.  Its canonical embedding $\delta_E:E\to\FBL[E]$ obeys
\[
 \delta_E(x)(\operatorname{Re}z^*)=z^*(x),\qquad x\in E,\ z^*\in E^*.
\]
For $T:E\to E$, the universal property produces $\overline T$ with
$\overline T\delta_E=\delta_ET$.  In the functional model,
\begin{equation}\label{eq:composition}
 (\overline T f)(\operatorname{Re}w^*)
 =f(\operatorname{Re}T^*w^*).
\end{equation}

The source proves inclusion of approximate point spectra, handles positive residual spectral values, and asks whether the full spectrum is always included.  It obtains only $|\lambda|\in\sigma(\overline T)$ for a general $\lambda\in\sigma(T)$.

\begin{center}
\includegraphics[width=0.72\textwidth]{figures/open_question_crop.png}
\end{center}

\section*{The Fourier-evaluation functional}

\begin{lemma}\label{lem:fourier}
Suppose $z^*\in E^*\setminus\{0\}$ and
$T^*z^*=\lambda z^*$.  Then $\lambda$ is an eigenvalue of $\overline T^*$.
\end{lemma}

\begin{proof}
Write $\lambda=re^{i\theta}$ with $r\geq0$; when $r=0$, choose $\theta=0$.  For $t\in\R$, let
\[
 \epsf_t(f)=f(\operatorname{Re}(e^{it}z^*)),
 \qquad f\in\FBL[E].
\]
These evaluations are uniformly bounded: $\|\epsf_t\|\leq\|z^*\|$.  Moreover $t\mapsto\epsf_t(f)$ is continuous, by the functional representation of $\FBL[E]$.  Hence
\begin{equation}\label{eq:Phi}
 \Phi(f)=\frac1{2\pi}\int_0^{2\pi}e^{-it}
 f(\operatorname{Re}(e^{it}z^*))\,dt
\end{equation}
defines a bounded linear functional.

It is nonzero.  Indeed, the canonical-evaluation identity gives, for every $x\in E$,
\[
 \Phi(\delta_E x)
 =\frac1{2\pi}\int_0^{2\pi}e^{-it}(e^{it}z^*)(x)\,dt
 =z^*(x).
\]
Choose $x$ with $z^*(x)\neq0$.

Using \eqref{eq:composition}, complex linearity of $T^*$, and positive homogeneity of every free-lattice function under the real scalar $r$, we obtain
\begin{align*}
 \Phi(\overline T f)
 &=\frac1{2\pi}\int_0^{2\pi}e^{-it}
 f(\operatorname{Re}T^*(e^{it}z^*))\,dt\\
 &=\frac r{2\pi}\int_0^{2\pi}e^{-it}
 f(\operatorname{Re}(e^{i(t+\theta)}z^*))\,dt\\
 &=re^{i\theta}\Phi(f)=\lambda\Phi(f).
\end{align*}
The last substitution is $s=t+\theta$, using periodicity.  Thus
$\overline T^*\Phi=\lambda\Phi$ with $\Phi\neq0$.
\end{proof}

\section*{Full spectral inclusion}

\begin{theorem}\label{thm:main}
For every complex Banach space $E$ and every bounded operator $T:E\to E$,
\[
 \sigma(T)\subseteq\sigma(\overline T).
\]
\end{theorem}

\begin{proof}
Fix $\lambda\in\sigma(T)$.  If $\lambda\in\sigma_a(T)$, take unit vectors $x_k$ with $(\lambda-T)x_k\to0$.  Since $\delta_E$ is an isometry and intertwines $T$ with $\overline T$,
\[
 \|(\lambda-\overline T)\delta_E x_k\|
 =\|\delta_E((\lambda-T)x_k)\|\longrightarrow0.
\]
Thus $\lambda\in\sigma_a(\overline T)\subseteq\sigma(\overline T)$.

Now suppose $\lambda\notin\sigma_a(T)$.  Then $\lambda-T$ is bounded below, so its range is closed.  It is not onto because $\lambda\in\sigma(T)$; hence its range is a proper closed subspace.  Hahn--Banach supplies $z^*\in E^*\setminus\{0\}$ annihilating $(\lambda-T)E$.  Equivalently,
$T^*z^*=\lambda z^*$.  Lemma~\ref{lem:fourier} gives
$\lambda\in\sigma_p(\overline T^*)$, which implies
$\lambda\in\sigma(\overline T)$.
\end{proof}

\section*{Why the averaging is the missing step}
A single evaluation at $\operatorname{Re}z^*$ detects $T^*z^*=\lambda z^*$ only when $\lambda\geq0$, because free-lattice functions are merely positively homogeneous on the real dual.  Rotating $z^*$ around its circle and taking the first Fourier coefficient projects precisely onto the complex-linear component: on $\delta_E(E)$ it equals $z^*$, while $T^*$ rotates the circle by $\arg\lambda$ and dilates it by $|\lambda|$.  This restores the full multiplier $\lambda$, rather than only its modulus.

\section*{Audit notes}
\begin{itemize}[leftmargin=1.4em,itemsep=2pt]
\item The integral in \eqref{eq:Phi} is a scalar integral of uniformly bounded continuous functions; it defines a bounded complex-linear functional.
\item The factor $e^{-it}$, not $e^{it}$, is essential: after shifting $t$ by $\theta$ it produces $e^{i\theta}$.
\item An eigenvalue of $\overline T^*$ is sufficient for membership in $\sigma(\overline T)$; no claim about the point spectrum of $\overline T$ is made.
\item The argument covers $\lambda=0$ as well.  Then $T^*z^*=0$ and the Fourier functional is annihilated by $\overline T^*$.
\item This packet does not address the separate real-scalar question whether nonisomorphic $E,F$ can have lattice-isomorphic $\operatorname{FBL}[E],\operatorname{FBL}[F]$.
\end{itemize}

\section*{Source}
D. de Hevia and P. Tradacete, \emph{Free complex Banach lattices}, arXiv:2207.08090 (2022).

\end{document}
